\section{K1: Schwellen der Erkennung}
\label{app:report_k1}
\subsection*{Kurzfassung}
\begin{itemize}
\item Gewählt: weekly / long\_\allowbreak{}term \% = 1,02 / 14,20 bei Anteil 0,20 (aktuelle Config 1,02 / 14,20)
\item Einzeltag-Schwelle = 2,3 bei 4,9 \% Alarmtagen (Ziel 5 \%, aktuelle Config 2,3)
\item 4 Resets auf validation, 2,03 je Jahr (Grenze höchstens 4)
\end{itemize}
\subsection*{Setup}
\begin{itemize}
\item \textbf{Lauf:} k1\_\allowbreak{}007
\item \textbf{Systeme:} Analyse, Streak, Einordnung nach Regel und Reset wie im Zyklus, ohne Prognose und ohne Sprachmodell
\item \textbf{Zeitraum:} 2014-04-02 bis 2016-03-19 (718 Tage, validation), durchgehend mit Reset, Historie davor aus train
\item \textbf{Regel:} weekly und long\_\allowbreak{}term halten ohne Reset an gleich vielen Tagen an (6 von 7 Tagen über der Schwelle). Gewählt wird der größte Anteil mit höchstens 4 Resets
\item \textbf{Einzeltag-Schwelle:} ein Wert für daily und weekday: der größte Score des Tages schlägt an dem Anteil der 718 Tage an, der am nächsten bei 5 \% liegt. Scores aus dem durchgehenden Lauf mit Reset und den gewählten Niveau-Schwellen
\item \textbf{Kontrolle:} Szenarien an 11 Zeitpunkten im Abstand von 60 Tagen, je 40 Tage davor und 60 danach. Nicht Teil der Wahl
\item \textbf{Commit:} 978003e
\end{itemize}

\subsection*{Schwellen je Anteil anhaltender Tage und Resets im durchgehenden Lauf}
\begin{tablang}{V{1.32} V{0.99} V{1.15} V{0.82} V{0.82} W{0.90}}{Anteil anhaltender Tage & Schwelle weekly & Schwelle long\_\allowbreak{}term [\%] & Resets & Resets je Jahr & gewählt}
0,0 & 2,5 & 25,5 & 0 & 0,0 &  \\
0,0 & 2,3 & 23,2 & 0 & 0,0 &  \\
0,0 & 2,1 & 22,4 & 0 & 0,0 &  \\
0,1 & 1,9 & 21,8 & 1 & 0,5 &  \\
0,1 & 1,7 & 21,2 & 1 & 0,5 &  \\
0,1 & 1,6 & 20,1 & 1 & 0,5 &  \\
0,1 & 1,5 & 19,2 & 1 & 0,5 &  \\
0,1 & 1,4 & 18,4 & 1 & 0,5 &  \\
0,1 & 1,4 & 17,3 & 2 & 1,0 &  \\
0,1 & 1,4 & 17,3 & 3 & 1,5 &  \\
0,1 & 1,3 & 16,7 & 3 & 1,5 &  \\
0,1 & 1,2 & 16,1 & 3 & 1,5 &  \\
0,1 & 1,2 & 15,8 & 3 & 1,5 &  \\
0,1 & 1,2 & 15,3 & 4 & 2,0 &  \\
0,2 & 1,1 & 15,1 & 4 & 2,0 &  \\
0,2 & 1,1 & 14,9 & 4 & 2,0 &  \\
0,2 & 1,0 & 14,7 & 4 & 2,0 &  \\
0,2 & 1,0 & 14,5 & 4 & 2,0 &  \\
0,2 & 1,0 & 14,2 & 4 & 2,0 & ja \\
0,2 & 1,0 & 13,7 & 5 & 2,5 &  \\
0,2 & 1,0 & 13,0 & 5 & 2,5 &  \\
0,2 & 0,9 & 12,9 & 5 & 2,5 &  \\
0,2 & 0,9 & 12,5 & 5 & 2,5 &  \\
0,2 & 0,9 & 12,0 & 6 & 3,1 &  \\
0,3 & 0,8 & 11,7 & 6 & 3,1 &  \\
0,3 & 0,8 & 11,3 & 6 & 3,1 &  \\
0,3 & 0,8 & 11,1 & 6 & 3,1 &  \\
0,3 & 0,7 & 10,9 & 6 & 3,1 &  \\
0,3 & 0,7 & 10,6 & 6 & 3,1 &  \\
\end{tablang}

\subsection*{Resets mit den gewählten Schwellen}
\begin{tablang}{W{1.00}}{Reset am}
19.06.2014 \\
13.10.2014 \\
13.12.2014 \\
02.08.2015 \\
\end{tablang}

\subsection*{Gewählte Einzeltag-Schwelle und Anteil der Tage mit Alarm im durchgehenden Lauf}
\begin{tablang}{V{1.20} V{0.80}}{Schwelle & Tage mit Alarm [\%]}
2,3 & 4,9 \\
\end{tablang}

\subsection*{Kontrolle mit eingebauten Szenarien an 11 Zeitpunkten (nicht Teil der Wahl)}
\begin{tablang}{W{0.84} W{1.87} W{0.84} V{0.48} W{0.96}}{Szenario & soll & Reset wegen des Szenarios & Tage bis zum Reset & Ereignistag gemeldet}
Kontrolllauf & still bleiben & 0/11 & - & 0/11 \\
scaled & erkennen (Zurücksetzen) & 6/11 & 18,5 & 3/11 \\
ramp\_\allowbreak{}28 & erkennen (Zurücksetzen) & 7/11 & 32,0 & 0/11 \\
ramp\_\allowbreak{}28\_\allowbreak{}down & erkennen (Zurücksetzen) & 7/11 & 31,0 & 0/11 \\
ramp\_\allowbreak{}56 & still bleiben & 5/11 & 34,0 & 0/11 \\
ramp\_\allowbreak{}56\_\allowbreak{}down & still bleiben & 5/11 & 35,0 & 0/11 \\
peak & Einzelanomalie am Ereignistag & 2/11 & 19,0 & 8/11 \\
level & Einzelanomalie am Ereignistag & 2/11 & 19,0 & 8/11 \\
\end{tablang}

\subsection*{Hinweise}
\begin{itemize}
\item Kontrolle, Reset wegen des Szenarios: so oft setzt das System zurück, früher als zur selben Zeit ohne Szenario. Tage bis zum Reset: Median ab dem Ereignistag. Ereignistag gemeldet: der Tag selbst ist auffällig, nur für Spitze und ganzen Tag gedacht
\item Gezählt wird jeder Reset im durchgehenden Lauf, auch an echten Wechseln (Urlaub, Saison). Die Grenze ist eine Setzung über die Blindzeit nach einem Reset
\item Der Anteil anhaltender Tage legt nur das Verhältnis der beiden Schwellen fest. Er wird ohne Reset auf 718 Tagen bestimmt, gezählt wird mit Reset
\item Kontrolle: Die Zeitpunkte liegen in den Kalibrierdaten, die Zahlen sind eine Ansicht, kein Beleg. Bei 11 Zeitpunkten verschiebt schon ein einzelner Fall die Anteile deutlich
\item Die gewählten Schwellen (weekly\_\allowbreak{}streak\_\allowbreak{}threshold, long\_\allowbreak{}term\_\allowbreak{}streak\_\allowbreak{}threshold\_\allowbreak{}pct, unusual\_\allowbreak{}threshold) sind in NumericalAnalysisConfig zu übernehmen, gerundet
\end{itemize}
